\documentclass[12pt,a4paper]{article}
\usepackage[utf8]{inputenc}
\usepackage[spanish]{babel}
\usepackage{geometry}
\usepackage{xcolor}
\usepackage{titlesec}
\usepackage{fancyhdr}
\usepackage{enumitem}

% Configuración de márgenes e identidad visual
\geometry{left=2.54cm,right=2.54cm,top=2.54cm,bottom=2.54cm}
\definecolor{rojoUMB}{HTML}{8B0000}

% Configuración de encabezados y pies de página
\pagestyle{fancy}
\fancyhf{}
\lhead{\footnotesize Guía de Uso Rápido - Cartilla Interactiva}
\rhead{\footnotesize BioGALF Home Health}
\lfoot{\footnotesize Soporte Tecnológico: Giovanni Lineros}
\rfoot{\thepage}

% Estilo de títulos
\titleformat{\section}{\color{rojoUMB}\normalfont\Large\bfseries}{\thesection}{1em}{}
\titleformat{\subsection}{\color{rojoUMB}\normalfont\large\bfseries}{\thesubsection}{1em}{}

\title{
    \vspace{-2cm}
    \textcolor{rojoUMB}{\textbf{INSTRUCTIVO DE USO RÁPIDO (FISIOTERAPIA)}}\\
    \large \textit{Cartilla Clínica Digital para el Manejo del Parkinson}
}
\author{}
\date{1 de abril de 2026}

\begin{document}

\maketitle
\vspace{-1.5cm} % Ajuste de espacio

\section{1. Inicio de Sesión y Registro del Paciente}
\begin{itemize}
    \item \textbf{Escudo de Acceso Profesional:} Al abrir la aplicación, la pantalla estará bloqueada por seguridad. Lea el aviso legal y pulse el botón \textit{Declaro ser profesional y acepto los términos} para habilitar el sistema.
    \item \textbf{Perfil del Paciente:} Antes de iniciar la evaluación, diligencie el nombre completo, el Estadio de Hoehn y Yahr (I-III) y la Evaluación Socio-Ambiental (red de apoyo y barreras en el hogar).
    \item \textit{Recomendación:} Llene este perfil al inicio de la sesión para que el registro clínico quede estructurado desde el primer minuto.
\end{itemize}

\section{2. Uso de la Batería Clínica (Cálculo Automático)}
Abra la prueba que necesite desplegando las opciones. \textbf{Usted no necesita calcular ni sumar nada manualmente.}

\subsection*{TUG Cognitivo (Cronómetro Integrado)}
\begin{itemize}
    \item Pulse \textbf{Start} al iniciar la prueba y \textbf{Stop} al finalizar (use \textbf{Reset} si requiere repetir).
    \item La aplicación mostrará el tiempo y la interpretación automática:
    \begin{itemize}
        \item Mayor de 12.6 segundos: \textit{Alerta de alto riesgo de caída}.
        \item Menor o igual a 12.6 segundos: \textit{Riesgo bajo}.
    \end{itemize}
\end{itemize}

\subsection*{Escala de Berg (BBS)}
\begin{itemize}
    \item Verá 14 ítems, cada uno con botones de calificación del 0 al 4.
    \item Seleccione un valor por ítem. El sistema sumará automáticamente el puntaje total sobre 56 y mostrará el nivel de riesgo en tiempo real.
\end{itemize}

\subsection*{FGA y UPDRS}
\begin{itemize}
    \item \textbf{FGA:} Seleccione el puntaje (0 a 3) en los 10 ítems. La PWA calcula el total y genera un resumen funcional de la marcha.
    \item \textbf{UPDRS:} Ingrese el resultado en el campo de texto libre o numérico según su criterio.
\end{itemize}

\section{3. Piloto de Tecnologías RV}
Esta sección es un apoyo complementario. Al hacer clic en el botón \textbf{Lanzar Simulación Piloto}, se abrirá una ventana emergente indicando que el \textit{"Entorno de Simulación RV en Construcción - BioGALF"} está en desarrollo. Cierre la ventana y continúe con la valoración.

\section{4. Guardado Final de la Valoración en la Nube}
\begin{itemize}
    \item \textbf{Regla de Integridad:} El botón \textit{Guardar Valoración} permanecerá bloqueado por defecto. Se habilitará únicamente cuando usted diligencie \textbf{al menos 2 pruebas clínicas} de la batería (Ej. TUG + Berg).
    \item Una vez habilitado, al guardar, todos los datos se envían a la nube de \textbf{BioGALF} para garantizar la trazabilidad clínica.
\end{itemize}

\vspace{0.5cm}
\noindent\fbox{%
    \parbox{\textwidth}{%
        \textbf{Nota de Responsabilidad Profesional:} La aplicación automatiza cálculos e interpreta rangos de riesgo basándose en evidencia científica; sin embargo, la validación final y la decisión terapéutica continúan siendo responsabilidad exclusiva del fisioterapeuta tratante.
    }%
}

\end{document}